\subsection{Silver labels and annotation status}\label{subsec:silver-labels-and-annotation-status}

The calibration and evaluation of the knowledge-extraction pipeline does not rely on gold labels annotated by a domain expert. The reference labels, used to calibrate and evaluate the pipeline, are silver labels \pinktodo{Zheng et al. 2023}: annotations generated by prompting a strong language model, Claude Opus 4.7, at temperature zero. The resulting labels therefore represent a strong-LLM-judge approximation of which clinical findings should be extracted from each text chunk; but should not be treated as ground truth.

\pinktodo{Add temperature inside the silver label cache, and regenerate silver labels by Opus for the ILP-selected 15 papers again with temperature zero.}

\pinktodo{The silver finding cases from Opus are not run on the 15 correct papers anyway, we have to regenerate the silver-findings.}

\pinktodo{Add the citation here.}

Silver labels are useful, as producing expert-curated annotations at the same scale would not be within the scope of this thesis. The silver-labels make it possible to calibrate the thresholds and evaluate the pipeline. In place of expert annotation, the thesis uses three silver-quality artifacts \pinktodo{(Zheng et al., 2023)}, each with a distinct role. First is a silver finding set, described in Subsection~\ref{subsec:knowledge-extraction-calibration-set}, which is used to calibrate thresholds, hyperparameters, and agreement-based cascade configuration. Second is the held-out evaluation finding set, described in Subsection~\ref{subsec:knowledge-extraction-eval-set}, which is generated the same way, but on a disjoint set of papers, and used to measure strict-F1 scores of the full knowledge-extraction pipeline. The third is the relation-classification dataset, described in Subsection~\ref{subsec:relation-classification-eval-set}, which consists of 300 labelled claim-pairs, and is used to evaluate the relation classification directly.

All silver findings and labelled claim pairs are generated by Claude Opus 4.7, at temperature zero. To produce the two finding sets, a silver-generation prompt was applied to each usable text chunk, as explained in Subsection~\ref{subsec:knowledge-extraction-eval-set}. The relation-classification set is generated from given labels: the claim pairs are generated, based on a supporting, contradicting, or unrelated label. These artifacts provide the data for calibration and evaluation of the pipeline components, without requiring expert annotations.

None of the three datasets are treated as definitive ground truth. Since the silver labels are generated by a language model, every calibrated threshold and every measurement inherits the biases of Claude Opus 4.7: the calibrated thresholds inherit them because they are tuned on the silver calibration set, and the pipeline strict-F1 score measurements and the relation-classification accuracy inherit them as they are scored against silver references. Overfitting, however, is not a concern for either evaluation: the knowledge-extraction and relation-classification evaluation datasets are both disjoint from the calibration set. The relation-classification set has another caveat; the synthetically generated claim-pairs in this set might have cleaner relations among each other, compared to real claim pairs on the corpus: supporting claims might be very similar, contradicting claims very sharply opposing, and unrelated claims more plainly unrelated. Therefore, generalization from 300 claim pairs to a corpus-scale output cannot be assumed, which will further be discussed in Chapter~\ref{chapter:Discussion}.
